\section{Positioning Against Prior Surveys}
\label{sec:related}
Prior surveys already provide substantial coverage of the design domain. The distinctive question for the present review is therefore not whether another model-family taxonomy can be drawn, but how evidence of \emph{post-generation editability} varies with the delivered representation and the actual evaluation protocol.

Huang \etal~\citep{huang2023survey} organize graphic design intelligence into representation, understanding, and generation, at both element and whole-design levels. Their overview illustration (Fig.~1) relates those capabilities to the creative process, while their task examples and dataset table support the taxonomy. Shi \etal~\citep{shi2023layoutsurvey} instead center layout representation, deep generative model families, task conditions, datasets, and metrics. Their Fig.~9 uses multiple independent taxonomic axes: model family, content dependence, and application setting.

Tian and G{\"u}nther~\citep{tian2024vectorsurvey} provide a mathematically grounded survey of smooth vector representations and the creation--rasterization--vectorization pipeline. Their emphasis on unified notation is particularly useful for separating vector editing affordances from bitmap reconstruction. Li \etal~\citep{li2024criticalsurvey} apply critical interpretive synthesis, beginning with a literature screening flow (Fig.~1), then illustrating the relationship among four AI-in-design paradigms and a workflow roadmap (Figs.~4--6). Their synthesis of 33 included studies has a different scope and methodological purpose from a model-centric survey.

Zou \etal~\citep{zou2025fragment} divide the field into perception and generation, with detailed subsections for visual objects, typography, aesthetics, and layout. Their Fig.~1 shows the overall AI graphic-design pipeline, Fig.~2 visualizes publication trends and task distributions, and Fig.~3 summarizes layout-analysis methods. The paper also compares preceding surveys and tabulates technologies. Malashenko \etal~\citep{malashenko2025svgreview} focus on LLM-driven SVG generation, editing, and understanding, connecting a task taxonomy to benchmarks and empirical evaluations.

\begin{table*}[t]
\centering
\caption{Survey designs inspected to guide this manuscript. Figure descriptions are based on author/publisher versions, not reproduced artworks.}
\label{tab:surveys}
\footnotesize
\begin{tabularx}{\textwidth}{@{}p{0.20\textwidth}p{0.33\textwidth}X@{}}
\toprule
\textbf{Survey} & \textbf{Main organization} & \textbf{Figure/table strategy and implication}\\
\midrule
Huang \etal~2023 \citep{huang2023survey} & Representation $\rightarrow$ understanding $\rightarrow$ generation; element/design level & Overview of design intelligence (Fig.~1); task examples; dataset comparison table.\\
Shi \etal~2023 \citep{shi2023layoutsurvey} & Layout representation, model taxonomy, applications, datasets and metrics & Multi-axis taxonomy (Fig.~9), structured comparisons and datasets.\\
Tian \& G{\"u}nther 2024 \citep{tian2024vectorsurvey} & Mathematical SVG/vector representation, creation, rasterization and vectorization & Representation/pipeline diagrams and common mathematical notation.\\
Li \etal~2024 \citep{li2024criticalsurvey} & Screening, qualitative coding, four design paradigms, synthesis and roadmap & Screening flow (Fig.~1), paradigm relations and workflow roadmap (Figs.~4--6).\\
Zou \etal~2025 \citep{zou2025fragment} & Background, perception, generation, multimodal trends and future work & System pipeline (Fig.~1), task/time statistics (Fig.~2), method taxonomy (Fig.~3).\\
Malashenko \etal~2025 \citep{malashenko2025svgreview} & SVG generation, editing and understanding, benchmarks, comparisons & Task-specific benchmark organization; evaluation summaries.\\
\midrule
\textbf{Ours (working)} & Task/representation $\times$ edit operation $\times$ evaluation evidence & Original R--E--V crosswalk (Fig.~\ref{fig:rev}), edit-evaluation protocol (Fig.~\ref{fig:protocol}), evidence-tier and measurement tables.\\
\bottomrule
\end{tabularx}
\end{table*}

\paragraph{Illustration design.}
Surveys use both \emph{orientation graphics} (an early system pipeline or taxonomy) and \emph{evidence graphics} (paper counts, screening, methods/datasets comparisons). We adopt an original orientation figure in Fig.~\ref{fig:rev} and a separate measurement figure in Fig.~\ref{fig:protocol}. Unlike a generic taxonomy, each explicitly distinguishes the claims supported by an output format, a demonstrated edit, and an evaluated operation. We also retain the comparison tables in Secs.~\ref{sec:representation} and \ref{sec:evaluation}. A publication-time histogram or a PRISMA flowchart is deliberately omitted until the two-snapshot corpus and its search/deduplication records can support reproducible denominators; invented screening counts would falsely suggest a completed systematic review.
